% Options for packages loaded elsewhere
\PassOptionsToPackage{unicode}{hyperref}
\PassOptionsToPackage{hyphens}{url}
\PassOptionsToPackage{dvipsnames,svgnames,x11names}{xcolor}
\PassOptionsToPackage{space}{xeCJK}
\documentclass[
]{article}
\usepackage{xcolor}
\usepackage[margin=2.5cm]{geometry}
\usepackage{amsmath,amssymb}
\setcounter{secnumdepth}{-\maxdimen} % remove section numbering
\usepackage{iftex}
\ifPDFTeX
  \usepackage[T1]{fontenc}
  \usepackage[utf8]{inputenc}
  \usepackage{textcomp} % provide euro and other symbols
\else % if luatex or xetex
  \usepackage{unicode-math} % this also loads fontspec
  \defaultfontfeatures{Scale=MatchLowercase}
  \defaultfontfeatures[\rmfamily]{Ligatures=TeX,Scale=1}
\fi
\usepackage{lmodern}
\ifPDFTeX\else
  % xetex/luatex font selection
  \setmainfont[]{DejaVu Sans}
  \setmonofont[]{DejaVu Sans Mono}
  \ifXeTeX
    \usepackage{xeCJK}
    \setCJKmainfont[]{Noto Sans CJK SC}
  \fi
  \ifLuaTeX
    \usepackage[]{luatexja-fontspec}
    \setmainjfont[]{Noto Sans CJK SC}
  \fi
\fi
% Use upquote if available, for straight quotes in verbatim environments
\IfFileExists{upquote.sty}{\usepackage{upquote}}{}
\IfFileExists{microtype.sty}{% use microtype if available
  \usepackage[]{microtype}
  \UseMicrotypeSet[protrusion]{basicmath} % disable protrusion for tt fonts
}{}
\makeatletter
\@ifundefined{KOMAClassName}{% if non-KOMA class
  \IfFileExists{parskip.sty}{%
    \usepackage{parskip}
  }{% else
    \setlength{\parindent}{0pt}
    \setlength{\parskip}{6pt plus 2pt minus 1pt}}
}{% if KOMA class
  \KOMAoptions{parskip=half}}
\makeatother
\usepackage{longtable,booktabs,array}
\usepackage{caption}
\captionsetup[table]{skip=6pt}
\newcounter{none} % for unnumbered tables
\usepackage{calc} % for calculating minipage widths
% Correct order of tables after \paragraph or \subparagraph
\usepackage{etoolbox}
\makeatletter
\patchcmd\longtable{\par}{\if@noskipsec\mbox{}\fi\par}{}{}
\makeatother
% Allow footnotes in longtable head/foot
\IfFileExists{footnotehyper.sty}{\usepackage{footnotehyper}}{\usepackage{footnote}}
\makesavenoteenv{longtable}
\setlength{\emergencystretch}{3em} % prevent overfull lines
\providecommand{\tightlist}{%
  \setlength{\itemsep}{0pt}\setlength{\parskip}{0pt}}
\usepackage{bookmark}
\IfFileExists{xurl.sty}{\usepackage{xurl}}{} % add URL line breaks if available
\urlstyle{same}
% fallback for those not using the hyperref driver hyperxmp:
\makeatletter
\@ifundefined{xmpquote}{\newcommand{\xmpquote}[1]{#1}}{}
\makeatother
\hypersetup{
  colorlinks=true,
  linkcolor={Maroon},
  filecolor={Maroon},
  citecolor={Blue},
  urlcolor={Blue},
  pdfcreator={LaTeX via pandoc}}

\author{}
\date{}

\begin{document}

\section{筑基篇课后习题
XIV：奇偶性的本质与素数的奇偶性观测}\label{ux7b51ux57faux7bc7ux8bfeux540eux4e60ux9898-xivux5947ux5076ux6027ux7684ux672cux8d28ux4e0eux7d20ux6570ux7684ux5947ux5076ux6027ux89c2ux6d4b}

\begin{quote}
\textbf{声明 (Declaration)}: 本文\textbf{不代表任何数学结论}
(non-mathematical-claim), 仅为基于 Lean 4
形式化的一种\textbf{实验观测报告} (experimental observation report)。

{\def\LTcaptype{none} % do not increment counter
\begin{longtable}[]{@{}
  >{\raggedright\arraybackslash}p{(\linewidth - 6\tabcolsep) * \real{0.2500}}
  >{\raggedright\arraybackslash}p{(\linewidth - 6\tabcolsep) * \real{0.2500}}
  >{\raggedright\arraybackslash}p{(\linewidth - 6\tabcolsep) * \real{0.2500}}
  >{\raggedright\arraybackslash}p{(\linewidth - 6\tabcolsep) * \real{0.2500}}@{}}
\toprule\noalign{}
\begin{minipage}[b]{\linewidth}\raggedright
习题
\end{minipage} & \begin{minipage}[b]{\linewidth}\raggedright
主题
\end{minipage} & \begin{minipage}[b]{\linewidth}\raggedright
Lean 形式化
\end{minipage} & \begin{minipage}[b]{\linewidth}\raggedright
DOI
\end{minipage} \\
\midrule\noalign{}
\endhead
\bottomrule\noalign{}
\endlastfoot
exercise\_i\_14 & 奇偶性的本质与素数观测（R171） & PatParityPrime.lean &
10.5281/zenodo.21917255 \\
\end{longtable}
}
\end{quote}

\textbf{Exercise XIV for the Foundation-Building Chapter: The Essence of
Parity and Observing Primes through Parity}

\begin{quote}
习题定位：筑基篇课后习题系列第十四题。观测奇偶性的本质，然后从这个角度观测素数，选对二者共享的基点。全部新定理只锚筑基篇定理（R085/R141/R145/R170）+
mathlib 数论基础（Prime.eq\_two\_or\_odd），不用外部引理。
\end{quote}

\emph{2026-08-13 · Internal research exercise · Lean 4 / mathlib v4.32.2
· 8 new theorems PROVED no sorry · 算器神魂论筑基篇课后习题 XIV }

\begin{center}\rule{0.5\linewidth}{0.5pt}\end{center}

\subsection{习题陈述}\label{ux4e60ux9898ux9648ux8ff0}

\begin{quote}
观测奇偶性的本质，然后从这个角度观测素数，有可能需要选对二者共享的基点。唯一论点：\textbf{奇偶性
= 模 2 折叠（2 槽环，偶 = 对称对还原，奇 = 对称对 +
1）；素数在奇偶性投影下坍缩（除 2 外全奇数）；★二者共享基点 = 2（模 2
折叠的周期 = 素数域的唯一例外 = 临界线圆上的点）。}
\end{quote}

\begin{center}\rule{0.5\linewidth}{0.5pt}\end{center}

\subsection{零、总论：为什么奇偶性是模 2
折叠}\label{ux96f6ux603bux8bbaux4e3aux4ec0ux4e48ux5947ux5076ux6027ux662fux6a21-2-ux6298ux53e0}

奇偶性的 pat 本质：

\begin{enumerate}
\def\labelenumi{\arabic{enumi}.}
\tightlist
\item
  \textbf{偶数 = 对称对还原}：2n = n + n（R170
  哥德巴赫对称对）------偶数 = 对称对关于 n 折叠还原到 0（R085 折叠类
  0）。
\item
  \textbf{奇数 = 对称对 + 1}：2n + 1 = n + (n+1)------奇数 =
  相邻对称对之和。
\item
  \textbf{奇偶性 = 模 2 投影}：n ↦ n mod 2 ∈ \{0, 1\}------2
  槽环（R141：单位根 n 槽环，n=2）的两个槽。
\end{enumerate}

\textbf{奇偶性的本质 = 模 2 折叠}------数轴折叠到两个类 \{偶,
奇\}（折叠类，R085）。

\begin{center}\rule{0.5\linewidth}{0.5pt}\end{center}

\subsection{一、奇偶性的本质（PROVED）}\label{ux4e00ux5947ux5076ux6027ux7684ux672cux8d28proved}

\textbf{★核心：偶数 = 对称对还原；奇数 = 对称对 + 1；奇偶性 = 模 2
投影（2 槽环）。}

\subsubsection{论证}\label{ux8bbaux8bc1}

\begin{enumerate}
\def\labelenumi{\arabic{enumi}.}
\tightlist
\item
  \textbf{偶数 = 对称对之和}（R170/R085）：2n = n + n------偶数 =
  对称对关于 n 折叠还原（哥德巴赫对称对结构）。
\item
  \textbf{奇数 = 对称对 + 1}：2n + 1 = n + (n + 1)------奇数 =
  相邻对称对之和（偏离折叠类 1）。
\item
  \textbf{奇偶性 = 模 2 投影}（R141 2 槽环）：Even n ⟹ n \% 2 = 0；Odd n
  ⟹ n \% 2 = 1------数轴折叠到 2 槽环的两个槽。
\end{enumerate}

\subsubsection{形式化（PatParityPrime.lean，4
定理）}\label{ux5f62ux5f0fux5316patparityprime.lean4-ux5b9aux7406}

\begin{itemize}
\tightlist
\item
  \texttt{even\_is\_symmetric\_pair}：2n = n + n------偶数 =
  对称对还原。
\item
  \texttt{odd\_is\_symmetric\_pair\_plus\_one}：2n + 1 = n +
  (n+1)------奇数 = 对称对 + 1。
\item
  \texttt{even\_two\_slot}：Even n ⟹ n \% 2 = 0------偶数的 2 槽环 0
  槽。
\item
  \texttt{odd\_two\_slot}：Odd n ⟹ n \% 2 = 1------奇数的 2 槽环 1 槽。
\end{itemize}

\textbf{验证}：0 sorry。锚定 ring / mathlib Even.mod\_two\_eq\_zero /
Odd.mod\_two\_eq\_one。

\begin{center}\rule{0.5\linewidth}{0.5pt}\end{center}

\subsection{二、从奇偶性观测素数：2
是唯一偶素数（PROVED）}\label{ux4e8cux4eceux5947ux5076ux6027ux89c2ux6d4bux7d20ux65702-ux662fux552fux4e00ux5076ux7d20ux6570proved}

\textbf{★核心：素数集合在奇偶性投影下坍缩------除 2
外全部落在奇数类（mathlib
Prime.eq\_two\_or\_odd）。奇偶性''看不出''素数分布（信息坍缩），但指出素数域的基点
= 2（唯一例外）。}

\subsubsection{论证}\label{ux8bbaux8bc1-1}

\begin{enumerate}
\def\labelenumi{\arabic{enumi}.}
\tightlist
\item
  \textbf{2 是唯一偶素数}（mathlib Prime.eq\_two\_or\_odd）：素数 p ⟹ p
  = 2 ∨ Odd p。
\item
  \textbf{素数奇偶性坍缩}：素数 p ≠ 2 ⟹ p 奇数------除 2 外所有素数 ≡ 1
  (mod 2)。
\item
  \textbf{观测结果}：奇偶性投影下素数集合 = \{2\} ∪
  \{奇数\}------几乎无区分度，但 \textbf{2
  是唯一例外点}（素数域的奇偶性基点）。
\end{enumerate}

\subsubsection{形式化（PatParityPrime.lean，2
定理）}\label{ux5f62ux5f0fux5316patparityprime.lean2-ux5b9aux7406}

\begin{itemize}
\tightlist
\item
  \texttt{two\_unique\_even\_prime}：素数 p ⟹ p = 2 ∨ Odd p------2
  是唯一偶素数。
\item
  \texttt{primes\_parity\_collapse}：素数 p ≠ 2 ⟹ Odd
  p------素数奇偶性坍缩。
\end{itemize}

\textbf{验证}：0 sorry。锚定 mathlib Prime.eq\_two\_or\_odd'。

\begin{center}\rule{0.5\linewidth}{0.5pt}\end{center}

\subsection{三、★共享基点：2（PROVED）}\label{ux4e09ux5171ux4eabux57faux70b92proved}

\textbf{★核心：奇偶性与素数共享基点 = 2------模 2 折叠的周期 =
素数域的唯一例外 = 临界线圆上的点。}

\subsubsection{论证}\label{ux8bbaux8bc1-2}

\begin{enumerate}
\def\labelenumi{\arabic{enumi}.}
\tightlist
\item
  \textbf{奇偶性侧}：2 槽环的周期 = 2（R141：模 2 折叠由 2
  定义）------奇偶性的结构常数是 2。
\item
  \textbf{素数侧}：2
  是唯一偶素数（two\_unique\_even\_prime）------素数的奇偶性例外点是 2。
\item
  \textbf{几何侧}：2 在临界线圆上（R145 critical\_circle\_points：‖2-1‖
  = 1）------2 是框架的几何点。
\item
  \textbf{★二者在 2 处交汇}：奇偶性的周期（模 2）= 素数域的唯一例外 =
  共享基点。
\end{enumerate}

\subsubsection{形式化（PatParityPrime.lean，1
定理）}\label{ux5f62ux5f0fux5316patparityprime.lean1-ux5b9aux7406}

\begin{itemize}
\tightlist
\item
  \texttt{two\_shared\_basepoint}：(∀ p, Prime p → p = 2 ∨ Odd p) ∧
  ‖2-1‖ = 1------★共享基点 2。
\end{itemize}

\textbf{验证}：0 sorry。锚定 two\_unique\_even\_prime + R145。

\begin{center}\rule{0.5\linewidth}{0.5pt}\end{center}

\subsection{四、全景（组合定理）}\label{ux56dbux5168ux666fux7ec4ux5408ux5b9aux7406}

\textbf{★核心：偶数 = 对称对还原 ∧ 奇数 = 对称对 + 1 ∧ 2
在临界线圆上------奇偶性（模 2 折叠）观测素数：素数集合坍缩到 \{2\} ∪
\{奇数\}，★共享基点 = 2。}

\subsubsection{形式化（PatParityPrime.lean，1
定理）}\label{ux5f62ux5f0fux5316patparityprime.lean1-ux5b9aux7406-1}

\begin{itemize}
\tightlist
\item
  \texttt{parity\_prime\_perspective}：2n = n+n ∧ 2n+1 = n+(n+1) ∧ ‖2-1‖
  = 1------全景。
\end{itemize}

\textbf{验证}：0 sorry。合取项分别锚 even\_is\_symmetric\_pair /
odd\_is\_symmetric\_pair\_plus\_one / R145。

\begin{center}\rule{0.5\linewidth}{0.5pt}\end{center}

\subsection{五、习题解答总结}\label{ux4e94ux4e60ux9898ux89e3ux7b54ux603bux7ed3}

{\def\LTcaptype{none} % do not increment counter
\begin{longtable}[]{@{}llll@{}}
\toprule\noalign{}
论点 & 内容 & 锚定 & 标签 \\
\midrule\noalign{}
\endhead
\bottomrule\noalign{}
\endlastfoot
偶数 = 对称对还原 & 2n = n + n & R170/R085 & PROVED \\
奇数 = 对称对 + 1 & 2n + 1 = n + (n+1) & R085 & PROVED \\
奇偶性 = 模 2 投影 & 2 槽环 \{0,1\} & R141 & PROVED \\
★2 = 唯一偶素数 & 素数奇偶性坍缩 & mathlib Prime & PROVED \\
★共享基点 2 & 模 2 周期 = 素数例外 = 临界线圆 & R145 & PROVED \\
\end{longtable}
}

\textbf{习题的回答（一句话）}：\textbf{奇偶性的本质 = 模 2 折叠（偶 =
对称对还原 2n = n+n，奇 = 对称对 + 1 2n+1 =
n+(n+1)）；从奇偶性观测素数：素数集合坍缩到 \{2\} ∪ \{奇数\}（除 2
外全奇数）；★二者共享基点 = 2------模 2
折叠的周期（奇偶性结构常数）恰好是素数域的唯一偶素数例外，且 2
落在临界线圆上（R145）。}

\begin{center}\rule{0.5\linewidth}{0.5pt}\end{center}

\subsection{六、相关对照}\label{ux516dux76f8ux5173ux5bf9ux7167}

{\def\LTcaptype{none} % do not increment counter
\begin{longtable}[]{@{}
  >{\raggedright\arraybackslash}p{(\linewidth - 4\tabcolsep) * \real{0.3333}}
  >{\raggedright\arraybackslash}p{(\linewidth - 4\tabcolsep) * \real{0.3333}}
  >{\raggedright\arraybackslash}p{(\linewidth - 4\tabcolsep) * \real{0.3333}}@{}}
\toprule\noalign{}
\begin{minipage}[b]{\linewidth}\raggedright
文献
\end{minipage} & \begin{minipage}[b]{\linewidth}\raggedright
对应内容
\end{minipage} & \begin{minipage}[b]{\linewidth}\raggedright
备注
\end{minipage} \\
\midrule\noalign{}
\endhead
\bottomrule\noalign{}
\endlastfoot
\textbf{欧几里得}（《几何原本》） & 素数定义/无穷性 & 素数的经典基础 \\
\textbf{R085/R141/R145/R170}（筑基篇） & 折叠类/2 槽环/临界线圆/对称对 &
本框架定理 \\
\textbf{mathlib} \texttt{Prime.eq\_two\_or\_odd} & 2 是唯一偶素数 &
mathlib 数论基础 \\
\end{longtable}
}

\begin{quote}
诚实边界：本习题交付奇偶性结构观测（模 2 折叠/对称对/共享基点
2），非素数分布理论；``共享基点 2''是结构观测结论（PROVED
的是三个独立事实 + 交汇点）。
\end{quote}

\begin{center}\rule{0.5\linewidth}{0.5pt}\end{center}

\emph{筑基篇课后习题 XIV · 2026-08-13 · Lean 4 / mathlib v4.32.2 · 8
theorems PROVED no sorry（PatParityPrime.lean）· 配套 Lean
形式化：formal/Formal/Toolkit/PatParityPrime.lean（快照见
../lean/PatParityPrime.lean）}

\end{document}
